% RESEARCH_DRAFT_SHA256: 89f744ceae1f5d269b9fd1b64f0fea6a3d801434d1685e4ffa25a7485f7f479d
% Layout and source-reference wrappers expanded for the web renderer.
\documentclass[11pt]{article}
\usepackage[margin=25mm]{geometry}
\usepackage{fontspec,unicode-math}
\setmainfont{Latin Modern Roman}
\setmathfont{Latin Modern Math}
\usepackage{amsmath,amssymb,mathtools,xurl,hyperref}
\hypersetup{colorlinks=true,urlcolor=blue,linkcolor=blue}
\setcounter{secnumdepth}{0}
\setlength{\parindent}{0pt}
\setlength{\parskip}{5pt}
\setlength{\emergencystretch}{3em}
\newcommand{\sref}[2]{\hyperlink{source:#1:#2}{[S#2]}}
\newcommand{\cataloguescope}[1]{\paragraph{Source record.}#1}
\newcommand{\floor}[1]{\left\lfloor#1\right\rfloor}
\begin{document}
\section*{Research-attempt provenance}
This standalone research attempt retains the complete original definition body below,
including its historical source and preparation notes. Those notes describe the
earlier catalogue preparation; the new attempt and its verification are identified
separately after that body. No original ID, catalogue status, or source statement is changed.

% UnsolvedMath ID 20001170; source catalogue code AIM-COMPUTATION-0008
\section{Exact support and trace-controlled ellipsoidal reachability}\label{20001170}
{\small UnsolvedMath ID 20001170 · Computer Science}
\subsection{Definitions and mathematical statement}

\paragraph{Shared notation.}Standard number systems are $\mathbb N,\mathbb Z,\mathbb Q,\mathbb R,\mathbb C$. Specialized notation and index ranges follow the cited source and the definitions below.


A reachable set consists of states that can be reached by admissible controls and disturbances over a stated time horizon. Hamilton--Jacobi--Bellman equations describe optimal-control value functions; Hamilton--Jacobi--Isaacs equations add competing controls in a differential game. An ellipsoid is $\{x:(x-c)^{\mathsf T}Q^{-1}(x-c)\leq1\}$ for $Q\succ0$, and ellipsoidal calculus computes inner or outer approximations of evolving sets.
\paragraph{Record scope and qualifications.}
The catalogue title is retained as metadata. Named certificates, contracts, bounds or frameworks in that title are not asserted as results of the original source extract.
The excerpt is a broad discussion with no fixed controlled dynamics or accuracy criterion. Its mention of ellipsoidal methods does not specify exact support or trace bounds. The supplied catalogue marks this record open; its present research status was not independently checked. Its supplied reconstruction is labelled unverified.
The source text below is rendered literally where necessary; damaged formulas and composite extracts are not silently repaired into a new mathematical claim.
The supplied catalogue marks this formulation as an unverified reconstruction; missing assumptions are not certified by the source link.
\paragraph{Statement recorded in the supplied source.}
A.5 Kurzhanski, Alexander

Among the central topics of modern control theory are problems of control synthesis for complex systems. These include problems of control under uncertainty and conflict as well as those under complex constraints. The solutions to such problems are well formalizable within Hamiltonian techniques and Dynamic Programming ideas, being as a rule reduced to the solution of HJB and HJBI partial differential equations or/and variational inequalities. An example where such methods are successfully applied are problems of forward and backward reachability for uncertain systems with further passage to control synthesis, safety verification, measurement feedback and related issues. The advent of HJB methods for continuous systems is due to the introduction of appropriate theories for generalized viscosity-type solutions and their equivalents. Another important topic is the solution to problems in dynamics and control for set-valued systems in terms of Hamiltonian formalism and Dynamic Programming. At present there is a sharp necessity for a breakthrough in numerical methods for solving equations and variational inequalities of the HJB-HJBI types motivated by problems of the above. There is also an interest in numerical methods for the comparison principle in HJB theory, which allows to calculate upper and lower bounds to exact solutions of HJB equations. Another perspective is the calculation of set-valued solutions to problems in evolution dynamics, estimation and control. Together with my colleagues we have developed an ellipsoidal calculus aimed at such problems and closely connected to HJB theory. The calculus is applicable to systems with original linear structure, but its methods allow effective calculations in high dimensions as well as computer animation in high dimensions through computer windows. For nonlinear systems some comparison methods were indicated. 8

The discussion of such methods as well as those for stochastic dynamics are within my primary interests at the AIM workshop. \href{https://aimath.org/WWN/optimalcontrol/optimalcontrol.pdf}{[S2]}
\paragraph{Source status.}
The supplied catalogue labels this source record open; that metadata was checked against the supplied snapshot on 2026-10-11. The label is preserved for provenance and is not an independent verification of present open status or mathematical correctness.
\subsection{Short English statement}
The participant discusses numerical control, reachability, comparison principles, and ellipsoidal set approximations.
\subsection{Sources}
\begin{itemize}
\item[S1] Ulam AI and the UnsolvedMath Contributors, UnsolvedMath: A Curated Collection of Open Mathematics Problems, record 20001170 (AIM-COMPUTATION-0008). Catalogue curation and metadata: CC BY 4.0. The source extract is retained; definitions, English summaries and qualifications are editorial additions.
\url{https://huggingface.co/datasets/ulamai/UnsolvedMath}.
\item[S2] AIM Workshop Problem Lists, source linked by catalogue record AIM-COMPUTATION-0008.
\url{https://aimath.org/WWN/optimalcontrol/optimalcontrol.pdf}.
\end{itemize}
\paragraph{Preparation and review.}
The original source and UnsolvedMath contributors retain attribution for the source extract. Codex (GPT-6; exact serving revision not exposed) assisted with definitions, English summaries, scope notes and formatting. No new human mathematical authorship or independently verified solution is asserted. Permanent source ID 20001170 and code AIM-COMPUTATION-0008 checked against the complete supplied registry and current repository. Distinct upstream records, including repeated questions and broad or damaged extracts, are retained for complete coverage of the source CS category. A repeated record is not a new independent conjecture. Complete mathematical-equivalence and current-literature review has not been performed. Preparation and final review follow contribution prompts C and D, with source gaps disclosed. Structural and rendering checks do not establish a new theorem or open status.
\label{end:20001170}

\section{Literature and context check}
The participant contribution in the primary AIM optimal-control workshop
 PDF was consulted. Its broad scope is retained above. The following explicitly
 chosen finite or smooth model tests one route suggested by the agenda; it is
 not a reconstruction of a missing theorem. The elementary arguments are not
 claimed novel, and no current-resolution audit of the whole field is asserted.
\section{Research attempt: Directionwise exact support does not give one exact ellipsoid}
For a positive semidefinite matrix $Q$, let $E(Q)=\{Q^{1/2}u:\|u\|\le1\}$,
including degenerate ellipsoids. Its support in direction $l$ is
$h_{E(Q)}(l)=\sqrt{l^TQl}$ by Cauchy--Schwarz, attained by a vector parallel
to $Q^{1/2}l$ when nonzero. The support of a Minkowski sum is the sum of
supports because its two maximizations are independent.

For $a=l^TQ_1l$, $b=l^TQ_2l$ and $\lambda>0$,
\[
 (\sqrt a+\sqrt b)^2\le(1+1/\lambda)a+(1+\lambda)b,
\]
since $a/\lambda+\lambda b-2\sqrt{ab}=(\sqrt{a/\lambda}-\sqrt{\lambda b})^2$.
Thus $E(Q_1)+E(Q_2)$ is contained in $E(Q(\lambda))$, where
$Q(\lambda)=(1+1/\lambda)Q_1+(1+\lambda)Q_2$.
For a direct containment justification, every point of the sum satisfies all
support inequalities of the outer ellipsoid. A point outside a closed convex
ellipsoid violates one: using its nearest point, the difference vector defines
a separating supporting hyperplane. The nearest-point first-order inequality
follows by differentiating squared distance along segments in the ellipsoid.

If both traces are positive, differentiating
$\operatorname{tr}Q(\lambda)$ shows its minimum occurs at
$\lambda=\sqrt{\operatorname{tr}Q_1/\operatorname{tr}Q_2}$; the derivative
changes from negative to positive there. But equality of support in a direction
requires $\lambda=\sqrt{a/b}$ when $a,b>0$, generally a different value.
Take $Q_1=\operatorname{diag}(1,0)$ and $Q_2=\operatorname{diag}(0,1)$.
Their sum set is the square $[-1,1]^2$, whereas the trace-minimizing outer
ellipsoid is the radius-$\sqrt2$ disk. It agrees with the square's support in
diagonal directions and strictly overestimates it on either axis.

Therefore assembling exact directional answers or optimizing one scalar shape
criterion does not produce a single exact ellipsoidal reachable set. This is
already visible for the linear control system whose two independent bounded
integrated inputs generate the two coordinate segments. General nonlinear
reachability requires additional error propagation not addressed by this sum.
\section{Verification and adversarial audit}
The support and outer-bound formulas were checked for zero directional
variance. Zero trace means $Q=0$ and is handled by omitting that summand, not
dividing by zero. The square/disk example verifies both equality and strict
inequality directions.
\paragraph{Audit corrections.} A trace-optimal outer ellipsoid was kept distinct from exact reachability or optimality for another error criterion.
\paragraph{External computation.} None; symbolic arguments and boundary checks only.
\section{Scope and remaining gaps}
The broad uncertain, competing-control and nonlinear HJB agenda is not
specified as one system. The first gap is a tractable approximation with an
explicit error criterion maintained through the actual dynamics; trace optimality
for one auxiliary sum is not that guarantee.
\paragraph{Integration review.} Codex (GPT-6; exact serving revision not exposed) reviewed the written arguments and expanded layout and source-reference wrappers. Repository metadata and the production web renderer were checked using Python 3.13 and JavaScript with MathJax 3.2.2. This model-assisted review is not independent human proof verification; no PDF compilation or formal proof checking is asserted.

\section{Completion Estimate}
\noindent\textbf{COMPLETION: 0\%}
The calculations below the preserved source statement investigate a concrete route,
but establish no substantiated new progress toward the entire original target.
Known special cases and failed sufficient conditions are not credited as a solution.
\section{References}
The original source references and catalogue attribution are preserved in the
Sources section above. The following URLs were actually consulted for this attempt;
retention of an additional original reference does not claim fresh consultation of it.
\begin{itemize}
\item \url{https://aimath.org/WWN/optimalcontrol/optimalcontrol.pdf} (consulted 2026-10-11). 
\end{itemize}
\section{Model and contributor disclosure}
OpenAI gpt-6-astra, requested reasoning configuration ultra, generated this research
attempt on 2026-10-11. The exact serving revision is not exposed. Codex GPT-6
performed orchestration, repository-format review, and a model-assisted mathematical
audit of the written arguments; its exact serving revision
is also not exposed. The mathematical route, boundary checks and final adversarial
self-audit followed contribution prompts A then B in
\texttt{docs/contribution\_prompts.md}, after reading \texttt{CONTRIBUTING.md}.
Python 3.14.2 assembled the standalone documents; any mathematical computation is
specified in Verification. No theorem prover was used. No human mathematical
authorship, substantial human proof correction, or independent human proof checking
is asserted. Model-assisted review, source retrieval, and structural checks do not
establish correctness or novelty. Full private reasoning transcripts are unavailable
as reproducibility artifacts; the complete arguments, inputs and outcomes actually
used are stated here or in the companion batch audit.
\section{Final status}
Unresolved as a new research contribution. This is a concrete route investigation,
not a claim of new complete resolution or a current-status certification. A known
answer reproduced for context does not become a new contribution. The catalogue's
source-status label is preserved; this attempt status does not assert that a
source-labelled solved question is open.
\end{document}
